\scanpage{150}
\[
\tag{57}
\begin{pmatrix}
 A+2\unsure{\nu}C&B+2\unsure{\nu}D\\
 C&D
\end{pmatrix}
=\varepsilon\begin{pmatrix}R&S\\T\,(=S)&U\end{pmatrix}.
\]

But this condition implies first that the matrix on the left-hand side is symmetric (since $S=T$) i.e.
\[
2\nu D=C-B
\]
a condition which is also equivalent, by (54), to
\[
\tag{58}
\inserted{2\nu D=1-D\quad\text{i.e.}}\qquad
\boxed{\mu D=1}\quad
\deleted{\mu D=\unsure{1}\;\illegible{1}\qquad\unsure{\mu}=\unsure{1}}\quad
\text{i.e.}\quad D=1/\mu,\quad\mu=1/D
\]
\inserted{where $\mu=2\nu+1$,}
which already implies \uline{$\mu$ and $D$ invertible}.

\deleted{On the other hand, we must have}
Symmetrically, if we write (57) in the form
\[
\varepsilon\begin{pmatrix}A&B\\C&D\end{pmatrix}
=\unsure{\varepsilon}\begin{pmatrix}R-2\nu T&S-2\nu U\\T&U\end{pmatrix},
\]
the condition $C+D-B=1$ becomes, for the entries of the matrix $\varphi$
\[
T+U-(S-2\nu U)=\varepsilon\ ;
\]
i.e.
\[
(2\nu+1)U=\varepsilon\quad\text{i.e.}
\]
\[
\tag{59}
\mu U=\varepsilon\quad\text{i.e.}\quad U=\varepsilon/\mu,
\quad\inserted{\mu=\varepsilon/U}
\]
which is also equivalent, taking (58) into account, to the relation
\[
\tag{60}
D=\varepsilon U\quad\text{or}\quad U=\varepsilon D
\]
which is none other than the equality of the entries in the bottom rows of the matrices (57).
\margin{149}

\scanpage{151}
\mspage{365}
We can write out the relations (57), as \deleted{\illegible{1}} four relations on the $9$ variables $a,b,c,d,r,s,t,u,\mu$ in addition to the relations (48), (49);
\[
\tag{61}
\left\{\begin{gathered}
\boxed{\begin{aligned}
D&=\varepsilon U\\
\mu D&=1\\
C&=\varepsilon T
\end{aligned}}
\quad\begin{aligned}
\deleted{\mu\unsure{U}=\unsure{\varepsilon}}{i.e.}\quad\deleted{\unsure{D=1/\mu}}\quad\mu U=\varepsilon\\
&
\end{aligned}\\
A=\varepsilon\bigl(R-\deleted{\unsure{2\nu}}\underbrace{(\mu-1)}_{2\nu}\underbrace{T}_{C}\bigr)
\end{gathered}\right.
\]
\margin{\uline{N.B.} this last one is a consequence of the others, taking $C+D-B=1$ and $AD-BC=1$ into account}
i.e. -- \deleted{or} \inserted{again}
\[
\tag{62}
\left\{\begin{gathered}
\boxed{\begin{aligned}
\mu(c^2+d^2+cd)&=\delta\\
\unsure{c}^2+d^2+cd&=\varepsilon\delta(t^2+u^2)\\
\uline{ac}+\uline{bd}+\uline{ad}-(c^2+d^2+cd)&=\delta\varepsilon(rt+su)
\end{aligned}}\\
a^2+b^2+ab-(ac+bd+bc)=\delta\varepsilon\bigl(r^2+s^2-(\mu-1)(rt+su)\bigr)
\end{gathered}\right.
\]
\margin{consequence of the preceding ones}

When the signs $\varepsilon,\delta$ are fixed, we find \inserted{with (48), (49)} a system of $6$ equations in $9$ variables $a,b,c,d,r,s,t,u,\mu$ -- we should be able to check that the scheme (over $\operatorname{Spec}\mathbb Z$) of solutions of these equations is of relative dimension $3$, and of relative dimension $1$ over the ``factors'' $(a,b,c,d)$ resp. $(r,s,t,u)$. We thus have an adelic point of this absolute scheme -- a point with coefficients in $\widehat{\mathbb Z}$\ldots
\margin{no, it is of total relative dimension $4$}

\uline{Theorem} The multiplier $\mu$ is determined as a function of $\unsure{\xi|\zeta}\inserted{=\begin{pmatrix}A&B\\C&D\end{pmatrix}}$ (considered as an element of $\operatorname{Sl}(2,\widehat{\mathbb Z})$), by the relation
\[
\mu D=1\qquad(58)
\]
\unsure{\textup{(+)}}
and it is determined \uline{up to sign} by $\varphi=\begin{pmatrix}R&S\\T&U\end{pmatrix}$ by the relation
\[
\mu U=\varepsilon\qquad(59)\qquad\text{where }\varepsilon\in\{\pm1\},
\]
\margin{150}

\scanpage{152}
and where $\varepsilon$ is the sign which occurs in the formulas
\[
\left\{\begin{aligned}
\unsure{\xi|\zeta}&=\varepsilon l_0^\nu\varphi\qquad(56)\\
g\rho(g)\rho^2(g)&=\varepsilon\qquad(45).\\
\inserted{g}&\inserted{=\varepsilon\alpha\rho(\alpha)^{-1}}
\end{aligned}\right.
\]
We have $\varepsilon=1$ iff $\mu\equiv1\ (4)$, $\varepsilon=-1$ iff $\mu\equiv-1\ (4)$.
\inserted{i.e. $\nu\equiv0\ (2)$}
\inserted{i.e. $\nu\equiv1\ (2)$}

\margin{in other words, we have $\mu\equiv\varepsilon\ (4)$ since $U=t^2+u^2\equiv1\ (4)$ hence $1/U\equiv(1)\ (4)$ hence $\mu=\varepsilon/U\equiv\varepsilon\ (4)$. This would be the time to include \unsure{also the|here the} relation $\mu\equiv\delta\ (3)$ cf. with the appropriate proof\ldots}

To prove this, we recall that
\[
U=t^2+u^2
\]
note that the only values modulo $4$ taken by $t^2$ and $u^2$ are $0$ and $1$, hence those taken by $t^2+u^2$ are $0,1$ and $2$, and the only unit among these is $1$. Thus \deleted{\unsure{$\mu$}} $1/\mu$ can be put in the form $\varepsilon(\unsure{t^2}+u^2)$ modulo $4$ iff $1/\mu\equiv\varepsilon\ (4)$ or again $\mu\equiv\varepsilon\ (4)$ q.e.d.

\deleted{Thus} We had seen in a quite analogous way, thanks to the formula $\mu D=1\quad(58)$ and
\[
D=\delta(c^2+d^2+cd)
\]
that $\mu\equiv\delta\ (3)$ i.e. $\delta=+1$ (resp. $-1$) iff $\mu\equiv1\ (3)$ (resp. $\mu\equiv-1\ \mathrm{mod}\ 3$). We \deleted{\unsure{could thus}} shall denote, for \unsure{all|the} prime numbers $p$, by $\unsure{\varepsilon_p|\beta_p}$ the value at $\mu$ of the ``quadratic character'' relative to $p$, equal to $1$ iff $\mu_p\in(\mathbb Z_p^*)^2$. We shall therefore have
\[
\tag{63}
\left\{\begin{aligned}
\delta&=ad-bc=\unsure{\varepsilon_3|\beta_3}(\mu)=\varepsilon_3\\
\varepsilon&=D/U=\unsure{\varepsilon_2|\beta_2}(\mu)=\varepsilon_2,
\end{aligned}\right.
\]
and knowing $\varepsilon_3$ (resp. $\varepsilon_2$, resp. both) is equivalent to knowing $\mu$ modulo $\unsure{3|2}$ (resp. modulo $4$, resp. modulo $12$).
\margin{$(\mathbb Z/12\mathbb Z)^*\simeq(\mathbb Z/4\mathbb Z)^*\times(\mathbb Z/3\mathbb Z)^*\simeq\{\unsure{\pm1}\}\times\{\pm1\}$}
\margin{151}

\scanpage{153}
\mspage{366}
We are also going to compute \deleted{\illegible{1}} (cf. formulas (45))
\[
\alpha=(l_0^\nu f)^{-1}\alpha_0=f^{-1}l_0^{-\nu}\alpha_0
\]
and
\[
g=\varepsilon\alpha\rho(\alpha)^{-1}
=\underbrace{\sigma_\infty(f)^{-1}}_{{}^t f}\,l_1^\nu\rho(f)
=(l_0^\nu f)^{-1}\unsure{\xi|\zeta}\,\rho(l_0^\nu f).
\]
We find
\[
\tag{64}
\begin{aligned}
\alpha&=\underbrace{\begin{pmatrix}u&-s\\-t&r\end{pmatrix}}_{f^{-1}}
\underbrace{\begin{pmatrix}1&2\nu\\0&1\end{pmatrix}
\begin{pmatrix}a&b\\c&d\end{pmatrix}}_{\begin{pmatrix}a+2\nu c&b+2\nu d\\c&d\end{pmatrix}}\\
&=\begin{pmatrix}u(a+2\nu c)-cs&u(b+2\nu d)-ds\\-t(a+2\nu c)+cr&-t(b+2\nu d)+dr\end{pmatrix},
\end{aligned}
\]
\margin{\uline{N.B.} Unlike $\alpha_0$ and $f$, and $\unsure{\xi|\zeta},\varphi$ which are deduced from them, $\alpha$ \deleted{\illegible{1}} shares with $g$ the advantage and the interest of being determined without ambiguity (up to sign) by an automorphism $u$ of $\widehat\Pi_{0,3}$, whereas $f$ etc. are \illegible{1} \unsure{choices of ambiguities} $f\mapsto l_0^n f$ etc.\ldots}
\[
\begin{aligned}
g&=\underbrace{\begin{pmatrix}r&t\\s&u\end{pmatrix}}_{\sigma_\infty(f^{-1})={}^t f}
\underbrace{\begin{pmatrix}1&0\\2\nu&1\end{pmatrix}}_{\unsure{l_1^\nu|l_0^\nu}}
\underbrace{\begin{pmatrix}t+u&-t\\\unsure{t+u-r-s}&\unsure{r-t}\end{pmatrix}}_{\rho(f)\ (\text{cf. \S43 formulas }(77))}\\
&=\begin{pmatrix}r&t\\s&u\end{pmatrix}
\begin{pmatrix}t+u&-t\\\unsure{\mu(t+u)-r-s}&-\mu t+r\end{pmatrix}\\
&=\begin{pmatrix}1+\mu t(t+u)&-\mu t^2\\-1+\mu u(t+u)&1-\mu ut\end{pmatrix},
\end{aligned}
\]
i.e.
\[
\tag{65}
g=\begin{pmatrix}1+\mu t(t+u)&-\mu t^2\\-1+\mu u(t+u)&1-\mu ut\end{pmatrix}.
\]

It appears that $g$ is expressed simply in terms of $f$ i.e. of $r,s,t,u$ directly, and not in terms of $\alpha_0$ or $\unsure{\xi|\zeta}$ i.e. via the $a,b,c,d$ or $A,B,C,D$. On the other hand, to express $\alpha$ (64), we must bring in both \inserted{the} $r,s,t,u$ and the $a,b,c,d$.

It was natural, in connection with the canonical homomorphism
\[
\mathcal N_{1,1}\longrightarrow\operatorname{Gl}(2,\widehat{\mathbb Z})
\]
(cf. \S39, (11)) to look for a matrix \deleted{\illegible{1}}, which by abuse of language I am going to denote by the same letter, $u$, as the automorphism under consideration of $\widehat\Pi_{0,3}$ with which we had started (just as I denote by the same letter the elements $g,f$ etc. of $\widehat\Pi_{0,3}$ or of $\unsure{\widehat{\mathfrak S}_{0,3}|\widehat{\mathcal S}_{0,3}}$, and their images in $\operatorname{Gl}(2,\widehat{\mathbb Z})'$, or even corresponding elements in $\operatorname{Gl}(2,\widehat{\mathbb Z})$ itself),
\margin{152}

\scanpage{154}
\[
u\in\operatorname{Gl}(2,\widehat{\mathbb Z})
\]
\deleted{\unsure{satisf}} with determinant equal to $\mu$, and inducing on $\operatorname{Sl}(2,\widehat{\mathbb Z})'$, by an inner automorphism, the same automorphism as \inserted{the} $u$ \inserted{at the outset} (by passage to the quotient), which is made explicit by the formulas
\[
\tag{66}\det u=\mu
\]
\[
\tag{67}u\sigma_\infty u^{-1}=\varepsilon'\sigma_\infty\qquad\varepsilon'\in\{\pm1\}
\]
\[
\tag{68}u\rho u^{-1}=\varepsilon''g\rho\qquad\varepsilon''\in\{\pm1\}
\]
In fact, \deleted{in} (68) \deleted{we \unsure{should have}} \inserted{is written}
\[
[u,\rho]=u\rho(u)^{-1}=\varepsilon''g
\]
and applying to both sides the operation $x\mapsto x\rho(x)\rho^2(x)$, and taking into account that $g\rho(g)\rho^2(g)=\varepsilon_2$, \deleted{we find} \inserted{and $x\rho(x)\rho^2(x)=1$ for $x=u\rho(u)^{-1}$,} we find
\[
\tag{68 bis}\varepsilon''=\varepsilon_2
\]
and (68) is thus equivalent to
\[
u\rho(u)^{-1}=\varepsilon_2g=\alpha\rho(\alpha)^{-1}
\]
\deleted{\unsure{whence}} i.e. to
\[
\tag{69}u=\alpha u_0\quad\text{with }[u_0,\rho]=1
\quad\bigl(\text{i.e. }u_0\rho(u_0)^{-1}=1\bigr).
\]

We have already remarked that (67), (68) determine $u$ only modulo multiplication by a scalar $\lambda\in\widehat{\mathbb Z}^*$, and taking (66) into account, $u$ is determined only modulo $\lambda$ satisfying $\lambda^2=1$ i.e. modulo ${}_2\widehat{\mathbb Z}^*\simeq\{\pm1\}^{\mathbb P}$ (where $\mathbb P$ is the set of prime numbers), but we nevertheless hope to hit upon a canonical choice.

In my first calculations, I hadn't yet \unsure{brought out|dug up} (69), where the condition $u_0$ commutes with $\rho$ was expressed in the form
\[
\tag{70}u_0=p\,\unsure{\mathrm{id}}+q\rho=\begin{pmatrix}p&q\\-q&p+q\end{pmatrix}\quad(p,q\in\widehat{\mathbb Z}),
\]
\margin{153}

\scanpage{155}
\mspage{367}

so that
\[
\det u_0=p^2+pq+q^2
\tag{71}
\]
and that (66) is written (taking (69) into account, \unsure{where}
$\det\unsure{\alpha}=\delta=\varepsilon_3$)
\[
p^2+pq+q^2=\mu\varepsilon_3
\tag{\unsure{72}}
\]
which, in view of the relation
\[
c^2+\unsure{cd}+d^2=\varepsilon_3/\mu
\qquad(52\text{ and }58)
\]
suggested relations of the form
\[
p=\deleted{\varepsilon}\mu c,\ q=\mu d,\quad
\text{or}\quad p=\mu d,\ q=\mu c.
\]
In fact, expressing (67) as an equation in $p,q$
(added \unsure{to|:} 72), \unsure{I found}
(up to sign of \unsure{course}
which here is inessential)
\[
p=\mu d,\ q=\mu c\qquad\text{i.e.}\qquad u_0=\mu(c\rho+d)
\tag{73}
\]
but working in the case $\varepsilon'=+1$ exclusively.
(I had thought I could conclude, a little too hastily, that the
sign $\varepsilon'$ was always $+1$, before realizing
that this was hardly compatible with class
field theory, and suspecting that we should have $\varepsilon'=+1$
iff $p\equiv1\ (4)$.) By making this rather cumbersome detour
via the $(a,b,c,d)$, I nevertheless arrived, after
many twists and turns (calculation errors etc.) at the
formula where there were no longer any $a,b,c,d$ at all,
and of a disconcerting simplicity:
\[
\uuline{u}=\mu(t\sigma_\infty+u)
=\mu\begin{pmatrix}u&t\\-t&u\end{pmatrix},
\tag{74}
\]
which manifestly corresponds to the case $\unsure{\varepsilon'_2|\varepsilon'}=1$, i.e.
the one where $\uline{u}$ \uline{commutes} with $\sigma_\infty$. As

\margin{($\Longleftrightarrow$ probably, $\varepsilon_2=1$)}
\margin{I underline the letter $u$ which denotes a matrix, to
distinguish it from the $u$ which denotes a scalar.}
\margin{154}

\scanpage{156}

\[
\det(t\sigma_\infty+u)=t^2+u^2=\varepsilon_2/\mu
\quad(=1/\mu\text{ if }\varepsilon_2=+1),
\]
we shall indeed have, in the case $\varepsilon_2=1$, \deleted{\unsure{therefore}}
\[
\det\uline{u}=\mu^2\det(t\sigma_\infty+u)=\mu^2\,1/\mu=\mu
\]
i.e., (66). On the other hand, we shall have
\[
\begin{aligned}
[u,\rho]&=[t\sigma_\infty+u,\rho]\\
&=
\underbrace{
\begin{pmatrix}u&t\\-t&u\end{pmatrix}
\begin{pmatrix}0&1\\-1&1\end{pmatrix}
}_{\begin{pmatrix}-t&u+t\\-u&u-t\end{pmatrix}}
\inserted{\mu}
\underbrace{
\begin{pmatrix}u&-t\\t&u\end{pmatrix}
\begin{pmatrix}1&-1\\1&0\end{pmatrix}
}_{\begin{pmatrix}u-t&-u\\u+t&-t\end{pmatrix}}.
\\
&=\mu\begin{pmatrix}
(u+t)^2-ut+t^2&-t^2\\
ut-t^2&u^2+t^2-ut
\end{pmatrix}
=\mu\begin{pmatrix}
\dfrac1\mu+\unsure{t(u+t)}&-t^2\\
t(u-t)&\dfrac1\mu-ut
\end{pmatrix}
\\
&=\begin{pmatrix}
1+\mu t(u+t)&-\mu t^2\\
\mu t(u\deleted{\unsure{t}}-t^{\deleted{\unsure{2}}})&1-\mu ut
\end{pmatrix},
\end{aligned}
\]
\margin{We could also \unsure{apply} the formulas §43 (78) ---}
and comparing with (65) we do indeed find
\[
[u,\rho]=g
\]
taking into account that
\[
\mu(u^2+t^2)=\varepsilon_2=1
\]
thus
\[
-1+\mu u(t+u)=\underbrace{-1+\mu u^2}_{-\mu t^2}+\mu ut
=\mu t(u-t)\ \unsure{-!|!}
\]
I would now like, without launching into
new calculations, to find ``by guesswork'' an expression
for $u$, in the \deleted{\unsure{case}} where $\varepsilon_2=-1$, and where we expect
to find in (67) the sign $\varepsilon'=-1$ --- i.e. that
$u$ be an \uline{anti}\unsure{similitude}\illegible{1}
\[
\uline{u}\,{}^t\uline{u}=-\det\uline{u}
\]
But \deleted{\unsure{and}} the matrix
\[
\tau_\infty=\begin{pmatrix}-1&0\\0&1\end{pmatrix}
\]
\margin{(it is \unsure{$=$} an \unsure{orthogonal} reflection)}
is the \unsure{typical case} of such an antisimilitude, and I
expect that the matrix
\[
\uline{u}=\mu\tau_\infty(t\sigma_\infty+u)
=\mu\begin{pmatrix}\unsure{1|-1}&0\\0&-1\end{pmatrix}
\begin{pmatrix}u&t\\-t&u\end{pmatrix}
=\mu\begin{pmatrix}\unsure{u|-u}&\unsure{t|-t}\\-t&-u\end{pmatrix}
\]
\margin{155}

\scanpage{157}
\mspage{368}

will do. We shall have, of course
\[
\begin{aligned}
u\,{}^tu
&=\mu^2\begin{pmatrix}u&t\\t&-u\end{pmatrix}
\begin{pmatrix}u&t\\t&-u\end{pmatrix}\\
&=\mu^2\begin{pmatrix}u^2+t^2&0\\0&u^2+t^2\end{pmatrix}
=\mu^2\begin{pmatrix}-1/\mu&0\\0&-1/\mu\end{pmatrix}
\end{aligned}
\]
i.e.
\[
u\,{}^tu=-\mu.\mathrm{Id}
\]
and
\[
\det u=\mu^2(-u^2-t^2)=\mu^2(1/\mu)=\mu
\]
taking into account that now
\[
u^2+t^2=\varepsilon_2/\mu=-\frac1\mu
\qquad\text{since }\varepsilon_2=-1.
\]
Finally to check (68) let us compute
\[
\begin{aligned}
[u,\rho]
&=\left[\begin{pmatrix}\unsure{u}&t\\t&-u\end{pmatrix},\rho\right]\\
&=\underbrace{
\begin{pmatrix}u&t\\t&-u\end{pmatrix}
\begin{pmatrix}0&1\\-1&1\end{pmatrix}
}_{\begin{pmatrix}-t&u+t\\u&-u+t\end{pmatrix}}
\unsure{\dfrac1\mu|\mu}
\underbrace{
\begin{pmatrix}-u&-t\\-t&u\end{pmatrix}
\begin{pmatrix}1&-1\\1&0\end{pmatrix}
}_{\begin{pmatrix}-u-t&\unsure{u|-u}\\u-t&t\end{pmatrix}}
\\
&=\mu\begin{pmatrix}
u(u+t)&t^2\\
\unsure{-(u^2+t^2)}-u(u-t)&u^2+t^2-ut
\end{pmatrix}
=\begin{pmatrix}
\mu u(u+t)&\mu t^2\\
1-\mu u(u-t)&-1-\mu ut
\end{pmatrix}
\end{aligned}
\]
(taking into account that $\mu(u^2+t^2)=-1$), \deleted{\unsure{whence|therefore}} do we indeed have
\[
[u,\rho]=-g
\]
taking
into account formula (65) which gives $g$ explicitly in
terms of $r,s,t,u$, we find that this
\unsure{expected|precise} formula is equivalent to
\[
2\mu ut=0\qquad\text{i.e.}\qquad ut=0
\]
really, this is screwy! But I see that if we
replace in the definition of $u$ $t$ by $-t$,
i.e. if we define
\[
\uline{u}=\mu(t\sigma_\infty+u)\tau_\infty
=\mu\begin{pmatrix}u&t\\-t&u\end{pmatrix}
\begin{pmatrix}1&0\\0&-1\end{pmatrix}
=\mu\begin{pmatrix}u&-t\\-t&-u\end{pmatrix}
\]
(instead of $\uline{u}=\mu\begin{pmatrix}u&t\\t&-u\end{pmatrix}$),
then $\uline{u}$ is \unsure{always}
an antisimilitude of determinant $\mu$, and
moreover now
\[
[u,\rho]=\begin{pmatrix}
\mu u(u-t)&\mu t^2\\
1-\mu u(u+t)&-1+\mu ut
\end{pmatrix}=-g
\]
\margin{156}

\scanpage{158}

which does indeed reduce to checking
\[
\mu u(u-t)=-(1+\mu t(t+u))
\]
i.e.
\[
\mu[u^2-\deleted{ut}+t^2+\deleted{tu}]=-1
\qquad\text{O.K.}
\]
In summary, we find

\uline{Theorem} Let $\alpha_2\in\mathbb Z/2\mathbb Z$ be defined by
$(-1)^{\alpha_2}=\varepsilon_2$
(thus $\alpha_2\equiv0\ (2)$ iff $\mu\equiv1\ (4)$,
$\alpha_2\equiv1\ (2)$ iff
$\mu\equiv-1\ (4)$). Put
\[
\uline{u}=\mu(t\sigma_\infty+u)\tau_\infty^{\alpha_2}
=\mu\begin{pmatrix}u&t\\-t&u\end{pmatrix}
\begin{pmatrix}1&0\\0&\varepsilon_2\end{pmatrix}^{\unsure{\alpha_2}},
\tag{75}
\]
\margin{$\displaystyle
=\begin{pmatrix}u&\varepsilon_2t\\-t&\varepsilon_2u\end{pmatrix}$}
\margin{\deleted{$\displaystyle
=\left\{\begin{array}{ll}
\mu\begin{pmatrix}u&t\\-t&u\end{pmatrix}&\alpha_2=0\\
\mu\begin{pmatrix}u&-t\\-t&-u\end{pmatrix}&\alpha_2=1
\end{array}\right.$}}
We then have the formulas
\[
\det\uline{u}=\mu
\tag{76}
\]
\[
[\uline{u},\sigma_\infty]=\unsure{-|\cdot}\varepsilon_2.\mathrm{id}
\qquad\text{i.e.}\quad\uline{u}(\sigma_\infty)=\varepsilon_2.\sigma_\infty
\tag{77}
\]
\[
[\uline{u},\rho]=\varepsilon_2.g
\qquad\text{i.e.}\quad\uline{u}(\rho)=(\varepsilon_2g).\rho
\tag{78}
\]

\uline{Corollary} Let $u$ be an automorphism \unsure{of the} group
with loops $\widehat{\Pi}_{0,3}$, \unsure{commuting} as outer actions with
$\unsure{\mathfrak S_3|\mathcal S_3}$ (i.e. with $\sigma_\infty$ and with $\rho$), $\mu$
its multiplier, let us also denote by
$u$ the canonical extension of $u$ to
$\unsure{\widehat{\mathfrak S}_{0,3}^{+}|\widehat{\mathcal S}_{0,3}^{+}}
\simeq\mathrm{Sl}(2,\mathbb Z)^{\wedge\prime}$,
inducing the identity on
$\unsure{\mathfrak S_3|\mathcal S_3}
\simeq\unsure{\widehat{\mathfrak S}_{0,3}^{+}|\widehat{\mathcal S}_{0,3}^{+}}/
\widehat{\Pi}_{0,3}$,
and consider the quotient $\mathrm{Sl}(2,\widehat{\mathbb Z})'$ of
$\mathrm{Sl}(2,\mathbb Z)^{\wedge\prime}$. Then the action of $u$ passes
to an action on the quotient $\mathrm{Sl}(2,\widehat{\mathbb Z})'$,
and the latter is induced by an
\deleted{\unsure{inner}}\inserted{``\unsure{inner|inner}''}
automorphism of $\mathrm{Gl}(2,\widehat{\mathbb Z})$, defined by a
\inserted{matrix} $\uline{u}\in\mathrm{Gl}(2,\widehat{\mathbb Z})$
such that $\det u=\mu$
\margin{157}

\scanpage{159}
\mspage{369}

(which describes the \uline{outer} action of $u$ on
$\mathrm{Sl}(2,\widehat{\mathbb Z})'$ in terms of the multiplier $\mu$).

\margin{\uline{N.B.} To check the \unsure{corollary}, we are
reduced to the case where $u$ commutes with $\sigma_\infty$,
if necessary conjugating by an inner automorphism \unsure{of}\illegible{1}\ldots}
\margin{We have done this, for the moment, for $u$
commuting with $\sigma_\infty$.}

In fact, we have much better still,
since we are going to define \uline{canonically}
\inserted{\uline{modulo sign}},
in terms of $u$, the matrix $\uline{u}$ by
formula (75), where \deleted{$t\text{ and }u\in$} $t,u\in\widehat{\mathbb Z}$ are
defined as follows: \deleted{\unsure{We have}} the datum of the automorphism $u$
defines $f\in\widehat{\Pi}_{0,3}$, modulo the operation $f\mapsto\ell_0^n f$
(for $n\in\widehat{\mathbb Z}$), and we again denote by $f$
the image of $f$ in $\mathrm{Sl}(2,\widehat{\mathbb Z})/\pm1$, which is thus
a matrix $\begin{pmatrix}r&s\\t&u\end{pmatrix}$, defined modulo sign,
We shall note that in $\mathrm{Sl}(2,\widehat{\mathbb Z})$ we have
\[
\ell_0^n\begin{pmatrix}r&s\\t&u\end{pmatrix}
=\begin{pmatrix}1&-2n\\0&1\end{pmatrix}
\begin{pmatrix}r&s\\t&u\end{pmatrix}
=\begin{pmatrix}r-2nt&s-2nu\\t&u\end{pmatrix}
\tag{79}
\]
($n\in\widehat{\mathbb Z}$), so $(t,u)$ does not change (modulo sign)
when we replace $f$ by $\ell_0^n f$ --- so $(t,u)$
is determined modulo sign by the automorphism $u$
of $\widehat{\Pi}_{0,3}$. \deleted{\illegible{1}} The analogous observation applies
moreover to $\alpha_0=\begin{pmatrix}a&b\\c&d\end{pmatrix}$, which is determined
modulo sign and modulo left multiplication
by an $\ell_0^n$ --- the \deleted{\illegible{1}} vector $(c,d)$ is
determined modulo sign, hence \unsure{the element}
\[
c\rho+d
\]
\margin{158}

\scanpage{160}

is determined modulo \uline{sign}, and \deleted{\illegible{1}}\inserted{unless I am mistaken}
we must have a relation
\[
\tag{80}
\uline{u}\overset{?}{=}\mu\alpha(c\rho+d)
=\mu f^{-1}\ell_0^{-\nu}\alpha_0(c\rho+d)
\]
(cf. (69) (73)), at least if $\varepsilon_2=1$ ---
otherwise we may perhaps have to correct by a factor $\tau_\infty$ somewhere.
\deleted{Now I have no \unsure{desire} for the moment to resume \unsure{my|my} \unsure{calculations|calculation} in this direction.}

\margin{N.B. this formula is \unsure{valid} \unsure{with} the corrective factor $\delta=\varepsilon_3$, cf. (82) below\ldots}

\[
\ell_0^{-\nu}\alpha_0(c\rho+d)
=
\begin{pmatrix}1&2\nu=\mu-1\\0&1\end{pmatrix}
\underbrace{
\begin{pmatrix}a&b\\c&d\end{pmatrix}
\begin{pmatrix}d&c\deleted{\unsure{+d}}\\-c&c+d\end{pmatrix}
}_{\displaystyle
\underbrace{
\begin{pmatrix}
ad-bc=\delta&ac+bc+bd=\delta B\\
0&c^2+cd+d^2=\delta D
\end{pmatrix}
}_{\displaystyle\delta\begin{pmatrix}1&B\\0&D\end{pmatrix}}
}
\]
\[
\ell_0^{-\nu}\alpha_0(c\rho+d)
=\delta\begin{pmatrix}1&B+2\nu D=C\\0&D\end{pmatrix}
\]
i.e.
\[
\tag{81}
\ell_0^{-\nu}\alpha_0(c\rho+d)
=\delta\begin{pmatrix}1&C\\0&D\end{pmatrix}
\]
whence
\[
\begin{aligned}
\mu\alpha(c\rho+d)
&=(\mu f^{-1})(\ell_0^{-\nu}\alpha_0(c\rho+d))
=\mu\delta\begin{pmatrix}u&-s\\-t&r\end{pmatrix}
\begin{pmatrix}1&C\\0&D\end{pmatrix}\\
&=\mu\delta\begin{pmatrix}u&uC-sD\\-t&-tC+rD\end{pmatrix}\\
&=\delta\begin{pmatrix}\mu u&-s+\mu uC\\-\mu t&r-\mu tC\end{pmatrix}
\end{aligned}
\]
Now (61)
\[
C=\varepsilon_2T=\varepsilon_2(rt+su)\qquad\text{whence}
\]
\[
\begin{aligned}
-s+\mu uC
&=-s+\varepsilon_2\unsure{\mu urt|\mu turt}+\varepsilon_2\mu u^2s\\
&\overset{\inserted{\unsure{\text{if }\varepsilon_2=+1}}}{=}
 s\bigl(\underbrace{-1+\unsure{\varepsilon_2u^2|\varepsilon_2\mu u^2}}_{-\varepsilon_2\mu t^2}\bigr)
 +\varepsilon_2\mu urt\\
&=\varepsilon_2\mu t\bigl(\overbrace{\unsure{-st+ur|st+ur}}^{+1}\bigr)\\
&=\varepsilon_2\mu t
\end{aligned}
\]

\margin{159}

\scanpage{161}
\mspage{370}

\[
\begin{aligned}
r-\mu tC
&\overset{\inserted{\unsure{\text{if }\varepsilon_2=1}}}{=}
\underbrace{r-\varepsilon_2\mu rt^2}_{r(1-\varepsilon_2\mu t^2)=\varepsilon_2r\mu u^2}
-\varepsilon_2\mu stu\\
&=\varepsilon_2\mu u(ru-st)=\varepsilon_2\mu u,
\end{aligned}
\]
thus
\[
\mu\alpha(c\rho+d)
=\delta\mu\begin{pmatrix}u&\varepsilon_2\unsure{t}\\-t&\varepsilon_2u\end{pmatrix}
\overset{\deleted{\unsure{\text{if }\varepsilon_2=1}}}{=}\delta\uline{u},
\]
i.e.
\[
\tag{82}
\uline{u}=(\unsure{\delta}\mu)\alpha(c\rho+d)
=(\delta\mu)(f^{-1}\ell_0^{-\nu}\alpha_0)(c\rho+d),
\]
\deleted{\unsure{if $\varepsilon_2=1$ (i.e. $\mu\equiv1\ (4)$)}}
\inserted{where $\delta=ad-bc\in\{\pm1\}$, $\delta=\varepsilon_3$
(thus $\delta\equiv\mu\ (3)$)}.

Formula which can also be written
\deleted{i.e.} with \unsure{account} \deleted{\illegible{1}} taken of (75)
\[
\tag{83}
\left\{\begin{aligned}
\deleted{\unsure{f}}(t\sigma_\infty+u)\inserted{\tau_\infty^{\alpha_2}}
&=\delta\alpha(c\rho+d)
\qquad\deleted{\unsure{\text{if }\varepsilon_2=1}}\\
\text{where}\quad\alpha&=f^{-1}\ell_0^{-\nu}\alpha_0,
\quad\deleted{\illegible{1}}\quad
\delta=\varepsilon_3\in\{\pm1\},\quad\delta=ad-bc,\quad\delta\equiv\mu\ (3)
\end{aligned}\right.
\]
or again
\[
\tag{84}
f(t\sigma_\infty+u)\tau_\infty^{\alpha_2}
=\delta\ell_0^{-\nu}\alpha_0(c\rho+d)
\]
\[
\begin{array}{c|c}
\underbrace{
\begin{pmatrix}r&s\\t&u\end{pmatrix}
\begin{pmatrix}u&t\\-t&u\end{pmatrix}
}_{\displaystyle\begin{pmatrix}1&rt+su\\0&t^2+u^2\end{pmatrix}}
\begin{pmatrix}1&0\\0&\varepsilon_2\end{pmatrix}
&
\underbrace{
\begin{pmatrix}1&2\nu\\0&1\end{pmatrix}
\begin{pmatrix}a&b\\c&\unsure{d}\end{pmatrix}
\begin{pmatrix}d&c\\-c&c+d\end{pmatrix}
}_{\displaystyle\begin{pmatrix}1&C\\0&D\end{pmatrix}}
\quad\text{cf. (81)}\\
\begin{pmatrix}1&rt+su\\0&t^2+u^2\end{pmatrix}\deleted{\begin{pmatrix}\unsure{u}&\unsure{t}\\-\unsure{t}&\unsure{u}\end{pmatrix}}&\\
\deleted{=\begin{pmatrix}1&T\\0&U\end{pmatrix}}&\\
\deleted{\unsure{=}}&\\
=\begin{pmatrix}1&\varepsilon_2T\\0&\varepsilon_2U\end{pmatrix}&
\end{array}
\]
and in this form the formula \unsure{was} in fact known,
as the essential content of the \deleted{formula} loop relation in matrix form (56) or (57)
\[
:\quad\varepsilon_2U=D,\qquad\varepsilon_2T=C\qquad!
\]

\uline{Remark} In fact, under a false appearance of arithmetic, we have done \unsure{by|for} the back door some scales in linear algebra, valid over

\margin{160}

\scanpage{162}

arbitrary rings and \unsure{formulable|formulable} probably under much more general and more ``\unsure{comprehensive}'' conditions --- concerning equations of the form
\[
[\alpha_0,\rho]=\xi
\]
where $\rho\in\unsure{\mathrm{Gl}}(2,A)$ is given, \deleted{\unsure{and}\illegible{1}}\inserted{as well as} $\xi$, and
$\alpha_0\in\mathrm{Gl}(2,A)$ is the unknown (possibly subject to conditions of the type $\det\alpha_0=\delta$, $\delta$ fixed), as well as simultaneous equations in $\alpha_0$ and $f$
\[
\tag{*}
[\alpha_0,\rho]=\varepsilon.\unsure{\mathcal{E}_0|\varepsilon_0}(\mu-1)[f,\sigma]
\qquad\text{cf. (56)}
\]
where $\rho,\sigma\in\mathrm{Gl}(2,A)$ are given, as well as a scalar $\varepsilon\in A^*$, and
$\lambda\longmapsto\unsure{\mathcal{E}_0|\varepsilon_0}(\lambda)$ is a ``one-parameter subgroup'', for example we can assume
\[
\unsure{\mathcal{E}_0|\varepsilon_0}(\lambda)=\begin{pmatrix}1&-\lambda\\0&1\end{pmatrix}=\varepsilon_0^\lambda,
\]
\[
\varepsilon_0=\begin{pmatrix}1&-1\\0&1\end{pmatrix}
\]
--- assuming if necessary a relation of the form $\sigma\rho=\varepsilon_0$.

Modulo some pleasant adjustments, all the formulas of the present \S\ should be able to be developed in this more general context
\inserted{(as well as substantial complements)}
(--- but I do not know whether I am going to spell out

\margin{161}

\scanpage{163}
\mspage{380}

here these \unsure{delightful} \unsure{scales} of linear algebra,
or even of scheme-theoretic interpretations of the kind that the scheme of solutions of (*)
(for $f,\alpha_0,\mu$ variable) is a smooth \inserted{relative} scheme of relative dimension 4, \unsure{which etc.}

I will only spell out, in the present adelic context, manifestly \unsure{too} special, the

\uline{Proposition} Let $\deleted{\unsure{a,b}}c,d\quad\deleted{\unsure{r,s}}t,u\in\widehat{\mathbb Z}$,
for them to \unsure{correspond} to data
$\alpha_0=\begin{pmatrix}a&b\\c&d\end{pmatrix}$,
$f=\begin{pmatrix}r&s\\t&u\end{pmatrix}$ satisfying
\[
\underbrace{ad-bc}_{\substack{\unsure{\text{definition}}\\=\varepsilon_3}}\in\{\pm1\},
\quad ru-st=1,
\quad\alpha_0\rho(\alpha_0)^{-1}=\varepsilon_2\ell_0^{\nu}(f\sigma_\infty(\unsure{f})^{-1})
\]
with $\varepsilon_2\in\{\pm1\}$ and $\unsure{\nu}\in\widehat{\mathbb Z}$
\unsure{suitable|suitable} (cf. (48) (49) (53)),
it \unsure{f|is necessary} and \deleted{\illegible{1}} it suffices that one have
\[
\tag{85}
c^2+\unsure{cd}+d^2=\pm(t^2+u^2)\quad\in\widehat{\mathbb Z}^*
\]
Then \deleted{\unsure{signs} $\varepsilon_2,\varepsilon_3$ and the scalar $\unsure{\nu}\in\widehat{\mathbb Z}$ are}
\inserted{$(\varepsilon_2,\varepsilon_3,2\nu+1)\in\{\pm1\}\times\{\pm1\}\times\widehat{\mathbb Z}^*$ is determined}
modulo sign (i.e. modulo passing to: $(-\varepsilon_2,-\varepsilon_3,-\nu)$)
\deleted{then uniquely determined as follows: One}
\inserted{by $a,b,c,d$, by the \deleted{condition} relation}
\inserted{(where $\mu=2\nu+1$)}
\deleted{defines \illegible{1}}
\[
\deleted{\varepsilon_3\unsure{\equiv|=}\unsure{(c^2+cd+d^2)}\unsure{\in\widehat{\mathbb Z}^*}}
\]
\[
\tag{86}
\frac1\mu=\varepsilon_2(t^2+u^2)=\varepsilon_3(c^2+cd+u^2)
\]
\margin{\unsure{by}}
\[
\deleted{\unsure{\mu|\frac1\mu}=\unsure{\varepsilon_3}(c^2+cd+d^2)=}
\]
\deleted{\illegible{3}}
\inserted{(which determines $\varepsilon_3,\mu$ when one has chosen $\varepsilon_2$)}
and one has
\[
\tag{87}
\left\{\begin{aligned}
\varepsilon_2&\equiv\mu\quad(4)\\
\varepsilon_3&\equiv\mu\quad(3)\\
\nu&=(\mu-1)/2
\end{aligned}\right.
\]

\margin{162}

\scanpage{164}

The solutions are then obtained as follows: one chooses
$\varepsilon_2,\varepsilon_3,\unsure{\mu|\mu_2}$ as above (two possibilities), \deleted{then} and arbitrarily $r,s\in\widehat{\mathbb Z}$ satisfying
\[
ur-ts=1
\]
(there are always some (for $r^2+s^2\in\widehat{\mathbb Z}^*$ --- hence
$\forall p\in\mathbb P$, $r_p$ and $s_p$ are not both in $p\mathbb Z_p$)
and if $r_0,s_0$ is a solution, the others are obtained in parametric form in $\lambda$
\[
\tag{87}
(r,s)=(r_0,s_0)\unsure{+|-}\lambda(\unsure{t|u},\unsure{u|t})
=(r_0\unsure{+|-}\lambda t,s_0\unsure{+|-}\lambda u).
\]
Having chosen $r,s$ hence $f=\begin{pmatrix}r&s\\t&u\end{pmatrix}\in\mathrm{Sl}(2,\widehat{\mathbb Z})$,

\deleted{
one deduces from this
\[
f\sigma_\infty(f)^{-1}=f\,{}^tf
=\begin{pmatrix}R&S\\T&U\end{pmatrix}
\]
where $R=\unsure{r}^2+s^2$, $S=T=rt+su$, $U=t^2+u^2$, and one successively determines the quantities
\[
\tag{88}
\left\{\begin{aligned}
D&=1/\mu=\varepsilon_3(c^2+cd+d^2)
=\deleted{\varepsilon_2\underbracket{\unsure{U}}_{\inserted{t^2+}}=\varepsilon_2(\unsure{t^2}+\illegible{1})}
\inserted{\varepsilon_2(t^2+u^2)=\varepsilon_2U}\\
C&=\varepsilon_2T\ (=\varepsilon_2S)\\
B&=C+D-1=\varepsilon_2(S+U)-1\\
A&=\frac1D(1+BC)=\frac1{\varepsilon_2U}(1+S(S+U)-\varepsilon_2S)
\end{aligned}\right.
\]
\margin{$\unsure{t^2+}$}
whence a matrix $\begin{pmatrix}A&B\\C&D\end{pmatrix}$ satisfying
$C+D-B=1$ and $D\in\widehat{\mathbb Z}^*$, and satisfying relation (57)
\[
\begin{pmatrix}A+2\nu C&B+2\nu D\\C&D\end{pmatrix}
=\varepsilon_2\begin{pmatrix}R&S\\T&U\end{pmatrix}
\]
which here reduces \unsure{to} the relation in the upper left corner
$A+2\nu C=\varepsilon_2R$ i.e.
}

\margin{163}

\scanpage{165}
\mspage{381}

it remains to choose $a,b$ so as to satisfy the two relations
$\det\begin{pmatrix}a&b\\c&d\end{pmatrix}=\varepsilon_3$
and $C=\varepsilon_2 T$ (\unsure{cf.} (61) (62))
\inserted{i.e.}
\[
\tag{88}
\left\{\begin{aligned}
d\uline{a}-c\uline{b}&=\varepsilon_3\\
(\unsure{c}+d)\uline{a}+d\uline{b}
&=\varepsilon_2\varepsilon_3
\underbrace{(\uline{r}t+\uline{s}u)}_{S}
+\underbrace{(c^2+cd+d^2)}_{1/\mu\varepsilon_3}
\end{aligned}\right.
\]
which is a system of two linear equations in $\uline{a},\uline{b}$,
whose determinant is as if by chance
$\unsure{c^2}+cd+d^2=1/\mu\varepsilon_3\in\widehat{\mathbb Z}^{*}$
--- which ensures a unique solution!

It seems more natural perhaps to start from the multiplier $\mu$, and to
look for the systems
$\begin{pmatrix}a&b\\c&d\end{pmatrix}$
$\begin{pmatrix}r&s\\t&u\end{pmatrix}$ corresponding to it,
but putting the emphasis on $(c,d)\ (t,u)$ in the first place.
Using $\mu$ one determines $\varepsilon_2,\varepsilon_3$ by
\[
\mu\equiv\varepsilon_2\ (2),\qquad
\mu\equiv\varepsilon_3\ (3),
\]
\deleted{\unsure{then}} which allows one to \deleted{\unsure{define}} represent
$1/\mu$ in the form
\[
1/\mu=\varepsilon_2(t^2+u^2)=\varepsilon_3(c^2+cd+d^2)
\]
(with $t,u,c,d\in\widehat{\mathbb Z}$, the systems $(t,u)$ as well as
$(c,d)$ ``\unsure{depend} on $1$ parameter'', \inserted{each}
in a very concrete sense that one could probably make precise if one cared to,
\deleted{\unsure{a wide choice}}\inserted{with}
explicit parameters \deleted{\illegible{1}} $\sigma,\tau$\ldots{} Once
$c,d,\unsure{t,u}$ have been chosen, one chooses $r,\unsure{s|t}$ as has

\margin{164}

\scanpage{166}

been said (with a parameter $\unsure{\lambda}$ \inserted{entering}
linearly), which then determines $r,s$ uniquely, in a way
that involves $\lambda$ \deleted{linearly} in a
linear-affine way.

In fact, for $c,d,t,u$ \inserted{and $\varepsilon_2,\varepsilon_3$}
fixed, \deleted{the equations}\inserted{so as to}
satisfy: $\varepsilon_2(t^2+u^2)=\varepsilon_3(c^2+cd+d^2)$,
the three equations relating $a,b,r,s$ are \unsure{linear-affine}
\[
\tag{89}
\left\{\begin{aligned}
d\uline{a}-c\uline{b}&=\varepsilon_3\\
u\uline{r}-t\uline{s}&=1\\
(c+d)\uline{a}+d\uline{b}
-\varepsilon_2\varepsilon_3(u\uline{s}+t\uline{r})&=c^2+cd+d^2
\end{aligned}\right.
\]
\margin{we have underlined the ``unknown coefficients''
$\uline{a},\uline{b},\uline{r},\uline{s}$ to be determined}
and the \deleted{determinant of this} system is ``of rank $3$''
\deleted{\unsure{det}}
--- the set of its solutions is an ``affine space of dimension $1$''
over $\widehat{\mathbb Z}$\ldots

\uline{Remark} We have not for the moment taken account of the congruence
conditions
\[
\begin{pmatrix}a&b\\c&d\end{pmatrix}\equiv1\ (2)
\qquad
\begin{pmatrix}r&s\\t&u\end{pmatrix}\equiv1\ (\unsure{2})
\]
which express that these matrices arise modulo signs by ``reduction''
from elements of
$\unsure{\widehat{\widetilde{\Pi}}_{0,3}}$, $\widehat{\Pi}_{0,3}$.
For the quantities $(c,d)$ \deleted{\unsure{fixed}} one finds the
consequences
\[
\left\{\begin{aligned}
d&\equiv\pm1\ (4)\quad\text{hence }d^2\equiv1\ (4)\\
c&\equiv0,2\ (4)\quad\text{hence }c^2\equiv0\ (4)
\end{aligned}\right.
\ \Longrightarrow\
\begin{aligned}
c^2+d^2+cd&\equiv1+cd\ (4)\\
&\equiv1+c\ (4)
\end{aligned}
\]
hence
\[
c^2+d^2+cd\equiv1+c\ \deleted{\unsure{(4)}}(4)
\quad\deleted{\unsure{\Longleftrightarrow}\ \unsure{cd}\ cd\equiv0\ (4)}
\quad\text{hence}
\]
hence
\[
\mu=\varepsilon_3(c^2+cd+d^2)\equiv
\deleted{\unsure{\Longleftrightarrow c\equiv0\ (4)}}
\varepsilon_3(1+c)\ (4)
\]
\deleted{(N.B. $d$ is invertible modulo $4$ since $d\equiv1\ (2)$),}

\margin{165}

\scanpage{167}
\mspage{382}

We thus see that knowledge of $c$ \deleted{\illegible{1}}
\inserted{modulo $4$}, and that of the sign $ad-bc=\varepsilon_3$,
allows one to find the sign $\varepsilon_2$ \deleted{\unsure{by}} by
\[
\deleted{\unsure{\varepsilon_2}\equiv\varepsilon_3\unsure{c}\quad
\varepsilon_2\equiv\varepsilon_3(1+c)\ (4)}
\]
\[
\tag{90}
\varepsilon_2\equiv\varepsilon_3(1+c)\ (4)
\quad\text{i.e.}\quad
\varepsilon_2\varepsilon_3\equiv1+c\ (4)
\]
If one sets out, for $c,d,t,u$ fixed, satisfying
\[
\tag{91}
\left\{\begin{aligned}
c&\equiv t\equiv0\ (2)\qquad d\equiv u\equiv1\ (2)\\
(c^2+\unsure{cd}+d^2)&\unsure{\equiv|=}\deleted{\unsure{(1+c}}
\varepsilon(t^2+u^2)\in\widehat{\mathbb Z}^{*}
\end{aligned}\right.
\]
where $\varepsilon$ \deleted{\illegible{1}}\inserted{$\in\{\pm1\}$}
is the sign defined by
\[
\varepsilon\equiv1+c\ (4)
\]
($\varepsilon=1$ if $c\equiv0\ (4)$, $\varepsilon=-1$ if
$c\equiv2\ (4)$), to find the corresponding $a,b,r,s$,
satisfying conditions (89), and moreover the congruence conditions
\[
\tag{92}
\unsure{b}\equiv\unsure{s}\equiv0\ (2)\qquad
a\equiv r\equiv1\ (2),
\]
one is led to lift an already given solution modulo $2$ of this system
of equations to a solution. In terms of the \inserted{$\widehat{\mathbb Z}$-}
torsor $\unsure{\mathcal T}$ of solutions, this means that one gives oneself
an element of the corresponding quotient $\mathbb Z/2\mathbb Z$-torsor
--- the set of its liftings is identified with a torsor under
$2\widehat{\mathbb Z}\simeq\widehat{\mathbb Z}$\ldots{} No mystery!

\margin{166}

\scanpage{168}

\uline{Remark} The additive group $\widehat{\mathbb Z}$ operates on the set
of solutions
$(\alpha_0,f)=\left(\begin{pmatrix}a&b\\c&d\end{pmatrix},
\begin{pmatrix}r&s\\t&u\end{pmatrix}\right)$
of the equations under consideration (48)(49)(56)
\[
\tag{93}
\det\alpha_0\in\{\pm1\},\quad\det f=1,\quad
\alpha_0\rho(\alpha_0)^{-1}
=\varepsilon\ell_0^{\nu}\bigl(f\sigma_\infty(f)^{-1}\bigr),
\quad\varepsilon\in\{\pm1\},
\]
by
\[
\tag{\unsure{93|94}}
\begin{aligned}
(\alpha_0,f)&\longmapsto(\alpha'_0,f')
=(\ell_0^n\alpha_0,\ell_0^nf)\quad\deleted{=(\unsure{a}-2nc,\ b-}\\
&=\left(\begin{pmatrix}a-2nc&b-2nd\\c&d\end{pmatrix},
\begin{pmatrix}r-2nt&s-2nu\\t&u\end{pmatrix}\right)
\end{aligned}
\]
($n\in\widehat{\mathbb Z}$).
Up to the factor $2$ in $2n$, this is precisely the indeterminacy in
\deleted{\illegible{1}} the set of solutions \inserted{of (93)}, when
$c,d,t,u$ are fixed (cf. (87) and the proposition that goes with it).
\inserted{(the factor $2$ is \unsure{ensured} \unsure{only when}
\unsure{one considers|in considering} solutions satisfying
$\alpha'_0\equiv f'_0\equiv\mathrm{Id}\ (2)$)}
One can \deleted{\unsure{thus}} say that what is determined
intrinsically \inserted{in
$\unsure{\mathrm{Gl}(2,\widehat{\mathbb Z})^{\pm1\prime}|\mathrm{Gl}(2,\widehat{\mathbb Z})^{\pm1}}$}
in terms of the \inserted{loop} automorphism $\uline{u}$ of
$\widehat{\Pi}_{0,3}$ commuting with $\sigma_\infty$, is neither more nor less
than \deleted{the \unsure{values}}
\[
\tag{95}
\begin{aligned}
\text{a)}\quad&\text{the multiplier }\mu\in\widehat{\mathbb Z}^{*}
\end{aligned}
\]
\deleted{\unsure{which gives} $\varepsilon_2=\unsure{\varepsilon_2|\xi_2}(\mu),\
\varepsilon_3=\unsure{\varepsilon_3|\xi_3}(\mu)$}
\inserted{(\unsure{hence|whence} $\unsure{\varepsilon_2}=\unsure{\varepsilon_2|\xi_2}(\mu),\
\varepsilon_3=\unsure{\varepsilon_3|\xi_3}(\mu)$)}

b) the quantities
\deleted{$\unsure{\delta}=\unsure{\varepsilon_3},\mu,\
c,d,t,u\in\widehat{\mathbb Z}$, \unsure{satisfying}}
the $c,d,t,u\in\widehat{\mathbb Z}$ satisfying the
everlasting relation (86)
\[
\tag{96}
\varepsilon_2(t^2+u^2)=\varepsilon_2(c^2+cd+d^2)=1/\mu
\]
and the congruence relations
\[
\tag{97}
c\equiv t\equiv0\ (2)\qquad d\equiv u\equiv1\ (2)
\]

It is true that we have seen that there is a \inserted{natural} way in
$\widehat{\Pi}_{0,3}$ to normalize the choice of the $f$, so as to remove
the \unsure{ambiguity|ambiguity} of the $\ell_0^nf$
(cf. \S\,41 \inserted{(24bis)} where this was touched upon --- I will come back to it
again later\ldots) but it seems to me \inserted{\uline{not very}}
interesting

\margin{167}
